\section{Anexos}

\subsection{Conexión de componentes al microcontrolador}
\begin{table}[H]
\centering
\caption{Conexión de componentes al XIAO ESP32C6}
\label{tab:conexion_xiao}
\begin{tabularx}{0.7\textwidth}{|X|X|}
\hline
\textbf{Componente} & \textbf{Pin del XIAO} \\
\hline
DS18B20 VCC  & 3.3V-OUT \\
\hline
DS18B20 DQ   & D0 \\
\hline
DS18B20 GND  & GND \\
\hline
PH-4502C VCC & 3.3V-OUT \\
\hline
PH-4502C GND & GND \\
\hline
PH-4502C PO  & D1 \\
\hline
MG-811 VCC   & 3.3V-OUT \\
\hline
MG-811 GND   & GND \\
\hline
MG-811 AOUT  & D2 \\
\hline
LED Azul     & D4 \\
\hline
LED Rojo     & D5 \\
\hline
Switch       & D6 \\
\hline
\end{tabularx}
\end{table}

\subsection{Reglas de lógica difusa para el pH y la temperatura}

\begin{enumerate}

% --- Muy Frío ---
\item IF Temp. es Muy Frío AND pH es Neutro THEN Condición es Inhibición (IN)
\item IF Temp. es Muy Frío AND pH es Inicial THEN Condición es Inhibición (IN)
\item IF Temp. es Muy Frío AND pH es Intermedio THEN Condición es Inhibición (IN)
\item IF Temp. es Muy Frío AND pH es Óptimo THEN Condición es Inhibición (IN)
\item IF Temp. es Muy Frío AND pH es Ácido THEN Condición es Riesgo Bajo (RIB)

% --- Frío ---
\item IF Temp. es Frío AND pH es Neutro THEN Condición es Deficiente (DEF)
\item IF Temp. es Frío AND pH es Inicial THEN Condición es Inicialmente Aceptable (IACE)
\item IF Temp. es Frío AND pH es Intermedio THEN Condición es Progresivamente Aceptable (PACE)
\item IF Temp. es Frío AND pH es Óptimo THEN Condición es Semi-Óptimo (SO)
\item IF Temp. es Frío AND pH es Ácido THEN Condición es Sobrefermentado (SOB)

% --- Ambiente ---
\item IF Temp. es Ambiente AND pH es Neutro THEN Condición es Deficiente (DEF)
\item IF Temp. es Ambiente AND pH es Inicial THEN Condición es Inicialmente Estable (IES)
\item IF Temp. es Ambiente AND pH es Intermedio THEN Condición es Progresivamente Estable (PES)
\item IF Temp. es Ambiente AND pH es Óptimo THEN Condición es Perfecto (PER)
\item IF Temp. es Ambiente AND pH es Ácido THEN Condición es Sobrefermentado (SOB)

% --- Caliente ---
\item IF Temp. es Caliente AND pH es Neutro THEN Condición es Deficiente (DEF)
\item IF Temp. es Caliente AND pH es Inicial THEN Condición es Inicialmente Aceptable (IACE)
\item IF Temp. es Caliente AND pH es Intermedio THEN Condición es Progresivamente Aceptable (PACE)
\item IF Temp. es Caliente AND pH es Óptimo THEN Condición es Semi-Óptimo (SO)
\item IF Temp. es Caliente AND pH es Ácido THEN Condición es Sobrefermentado (SOB)

% --- Muy Caliente ---
\item IF Temp. es Muy Caliente AND pH es Neutro THEN Condición es Peligro (PE)
\item IF Temp. es Muy Caliente AND pH es Inicial THEN Condición es Deficiente (DEF)
\item IF Temp. es Muy Caliente AND pH es Intermedio THEN Condición es Deficiente (DEF)
\item IF Temp. es Muy Caliente AND pH es Óptimo THEN Condición es Deficiente (DEF)
\item IF Temp. es Muy Caliente AND pH es Ácido THEN Condición es Peligro (PE)

\end{enumerate}


\subsection{Reglas de lógica difusa para el pH y el $CO_{2}$}

\begin{enumerate}

% --- CO2 Muy Bajo ---
\item IF CO2 es Muy Bajo AND pH es Neutro THEN Condición es Inactivo (INA)
\item IF CO2 es Muy Bajo AND pH es Inicial THEN Condición es Inactivo (INA)
\item IF CO2 es Muy Bajo AND pH es Intermedio THEN Condición es Inactivo (INA)
\item IF CO2 es Muy Bajo AND pH es Óptimo THEN Condición es Inactivo (INA)
\item IF CO2 es Muy Bajo AND pH es Ácido THEN Condición es Inestable (INE)

% --- CO2 Bajo ---
\item IF CO2 es Bajo AND pH es Neutro THEN Condición es Inactivo (INA)
\item IF CO2 es Bajo AND pH es Inicial THEN Condición es Saludable (SAL)
\item IF CO2 es Bajo AND pH es Intermedio THEN Condición es Saludable (SAL)
\item IF CO2 es Bajo AND pH es Óptimo THEN Condición es Saludable (SAL)
\item IF CO2 es Bajo AND pH es Ácido THEN Condición es Excesivo (EXC)

% --- CO2 Media ---
\item IF CO2 es Media AND pH es Neutro THEN Condición es Inestable (INE)
\item IF CO2 es Media AND pH es Inicial THEN Condición es Saludable (SAL)
\item IF CO2 es Media AND pH es Intermedio THEN Condición es Saludable (SAL)
\item IF CO2 es Media AND pH es Óptimo THEN Condición es Óptimo (OPT)
\item IF CO2 es Media AND pH es Ácido THEN Condición es Excesivo (EXC)

% --- CO2 Alto ---
\item IF CO2 es Alto AND pH es Neutro THEN Condición es Riesgo (RIE)
\item IF CO2 es Alto AND pH es Inicial THEN Condición es Saludable (SAL)
\item IF CO2 es Alto AND pH es Intermedio THEN Condición es Saludable (SAL)
\item IF CO2 es Alto AND pH es Óptimo THEN Condición es Óptimo (OPT)
\item IF CO2 es Alto AND pH es Ácido THEN Condición es Excesivo (EXC)

% --- CO2 Excesivo ---
\item IF CO2 es Excesivo AND pH es Neutro THEN Condición es Riesgo (RIE)
\item IF CO2 es Excesivo AND pH es Inicial THEN Condición es Excesivo (EXC)
\item IF CO2 es Excesivo AND pH es Intermedio THEN Condición es Excesivo (EXC)
\item IF CO2 es Excesivo AND pH es Óptimo THEN Condición es Excesivo (EXC)
\item IF CO2 es Excesivo AND pH es Ácido THEN Condición es Riesgo (RIE)

\end{enumerate}

\subsection{Descripción de las categorías de la fermentación del kéfir}

\begin{table}[htbp]
    \centering
    \caption{Explicación de las Etiquetas Resultantes}
    \label{tab:explicacion_etiquetas_globales}
    \begin{tabularx}{\textwidth}{|c|c|>{\raggedright\arraybackslash}p{3.8cm}|X|}
        \hline
        \textbf{Etiqueta} & \textbf{Abrev.} & \textbf{Significado Técnico Global} & \textbf{Implicación y Toma de Decisiones} \\
        \hline
        \textbf{Perfecto} & \textbf{PER} & Sincronía ideal entre cinéticas de acidificación, perfil térmico y desgasificación de $\text{CO}_2$. & Cosechar de inmediato para garantizar las máximas propiedades organolépticas y probióticas. \\
        \hline
        \textbf{Óptimo} & \textbf{OPT} & Desarrollo metabólico sobresaliente dentro de los márgenes óptimos de fermentación. & Cosecha próxima. Mantener el monitoreo constante sin alterar variables ambientales. \\
        \hline
        \textbf{Aceptable} & \textbf{ACE} & Proceso funcional estable con ligeras variaciones de tasa sin afectar la viabilidad del producto. & Continuar la fermentación. Permite el consumo directo sin requerir ajustes inmediatos. \\
        \hline
        \textbf{Incompleto} & \textbf{INC} & Desfasamiento por baja actividad microbiana o acidez insuficiente en relación con el tiempo. & Extender el tiempo de fermentación o ajustar ligeramente la temperatura para activar el cultivo. \\
        \hline
        \textbf{Caduco} & \textbf{CAD} & Inactividad metabólica prolongada o desintegración de las condiciones operativas de inocuidad. & El lote no alcanzó la acidificación adecuada. Se sugiere descartar el producto para evitar riesgos sanitarios. \\
        \hline
        \textbf{Riesgo} & \textbf{RIE} & Desequilibrio significativo por sobrefermentación, estrés térmico o hipercarbonatación. & Intervenir inmediatamente. Separar los granos de kéfir o refrigerar para frenar la cinética bacteriana. \\
        \hline
        \textbf{Peligro} & \textbf{PEL} & Desviación crítica irreversible con posible degradación hacia fermentación acética o contaminación. & Detener el proceso por completo, desechar el lote y esterilizar el contenedor e instrumental. \\
        \hline
    \end{tabularx}
\end{table}